% Exploratory results only. Requires booktabs and graphicx for the tables.
% Replace this section with the frozen multi-seed held-out study when available.
\subsection{Exploratory Benchmark Results}
We evaluated Tri-Fair-v8.3, the corrected Tri-Fair hybrid (implementation
version 8.4.1), and NSGA-II-PO-Fair on Bias-in-Bios and WinoBias using
Qwen3-30B and seed 72. Each method received a downstream-token ceiling of
one million tokens, comprising a 650,000-token search allowance and a
350,000-token reserve for completing a frozen final shortlist. All six
processes completed successfully. This screening experiment uses search-set
outcomes and does not constitute the final held-out comparison.

\paragraph{Comparison of approximation sets.}
Following the distinction between approximation-set quality and representative
prompt performance in MO-CAPO~\cite{bussing2026mocapo}, we report both search
hypervolume and in-sample R2 alongside the objectives of a single balanced
representative. The three loss coordinates are $(1-q,c/C_{\max},f)$, where
$q$ denotes macro-F1 on Bias-in-Bios and accuracy on WinoBias, $c$ denotes
configured inference cost, and $f$ denotes unfairness. Cost bounds are fixed
at 120 and 100, respectively, and the hypervolume reference point is
$(1.1,1.1,1.1)$. R2 uses 4,096 common Dirichlet preference vectors with seed
2026. These settings extend the original two-objective presentation to three
objectives and should not be interpreted as numerical replication of its
normalization protocol.

\begin{table*}[t]
\centering
\caption{Exploratory search results for seed 72. Each objective row refers to
the same balanced prompt. The indicators describe the two audited final
incumbents per method. No across-seed uncertainty is estimated.}
\resizebox{\textwidth}{!}{%
\begin{tabular}{llrrrrrr}
\toprule
Dataset & Method & Accuracy (\%) $\uparrow$ & Macro-F1 (\%) $\uparrow$ & Cost $\downarrow$ & Unfairness $\downarrow$ & Search HV $\uparrow$ & Search R2 $\downarrow$\\
\midrule
Bias-in-Bios & Tri-Fair-v8.3 & 72.54 & 72.97 & 18.8429 & 0.1719 & 0.767566 & 0.113761\\
 & Hybrid-v8.4.1 & 74.11 & 74.32 & 18.6238 & 0.1543 & 0.754139 & 0.118626\\
 & NSGA-II-PO-Fair & 74.33 & 74.55 & 18.9767 & 0.1981 & 0.719013 & 0.126262\\
\midrule
WinoBias & Tri-Fair-v8.3 & 86.72 & -- & 14.4830 & 0.2102 & 0.830346 & 0.098913\\
 & Hybrid-v8.4.1 & 85.35 & -- & 17.4530 & 0.1217 & 0.921181 & 0.079516\\
 & NSGA-II-PO-Fair & 86.33 & -- & 17.4530 & 0.1335 & 0.921952 & 0.079420\\
\bottomrule
\end{tabular}}
\end{table*}

On Bias-in-Bios, v8.3 obtained the strongest audited approximation set by
both search indicators. The hybrid's balanced representative nevertheless
improved all three objective values relative to the balanced v8.3 prompt.
Against NSGA-II-PO-Fair, its macro-F1 was 0.23 percentage points lower,
while its cost and unfairness were lower. This illustrates why conclusions
about one representative cannot be substituted for conclusions about an
entire approximation set.

On WinoBias, the hybrid and NSGA-II-PO-Fair yielded closely matched search
indicators, with NSGA-II-PO-Fair having slightly higher HV and lower R2.
Their balanced representatives had equal configured inference cost; the
hybrid had lower unfairness but accuracy lower by 0.98 percentage points.
The v8.3 representative offered lower cost and higher accuracy, accompanied
by higher unfairness. No method dominated across all objectives.

\paragraph{Quality-first representatives.}
MO-CAPO separately evaluates a maximum-quality prompt selected using
development performance. Applying that policy to the present search
candidates, v8.3 achieved 78.79\% accuracy and 78.83\% macro-F1 on
Bias-in-Bios, at cost 39.1957 and unfairness 0.1866. The corresponding
hybrid and NSGA-II macro-F1 values were 74.32\% and 74.55\%.
The quality-first v8.3 prompt is more expensive than its balanced
representative. Each policy therefore must be stated explicitly; objective
extrema from different prompts must not be combined into one row.

\paragraph{Search progress and limitations.}
The corrected hybrid completed nine logged steps on Bias-in-Bios and eleven
on WinoBias, using 576,414 and 637,595 downstream tokens. It now continues
past rejected proposals and terminates when a nonempty full candidate cannot
fit within the remaining search allowance. The remaining allowance and
unused audit reserve are reported rather than treated as consumed tokens.
The supplied archive contains no held-out predictions. Consequently, held-out
noisy R2, optimistic and pessimistic HV, and their approximation gap cannot
yet be reported. No significance claim is supported by this single-seed
screening experiment. The next stage evaluates the frozen final candidates
on validation data, followed by a prespecified multi-seed study and final
test evaluation after configuration freeze.
